This assignment taught me that using an AI coding agent well is not just about asking it to write code. At the start, I thought the main challenge would be getting the agent to produce something that actually ran. But I quickly realized the harder part was giving it enough structure to keep the project on track. I used AGENTS.md to set the long-term rules: scientific integrity, reproducibility, mathematical rigor, experiment requirements, and report quality. Then I used smaller, task-specific Markdown prompts to break the work into stages: finding the algorithm, implementing it, running experiments, reviewing the results critically, and preparing the final report. That separation mattered. The permanent rules told the agent how the project should always behave, while the task prompts told it what to do next. I also learned that vague instructions lead to vague results. Clear requirements, constraints, and acceptance criteria made a real difference.

I also saw how specialized AI skills can raise the quality of the work. Instead of depending only on the model\textquoteright{}s general knowledge, I gave it skills for literature research, machine learning, PyTorch implementation, experimental design, statistical analysis, scientific visualization, and academic writing. The prompts explained what needed to be done, while the skills gave guidance on how to do it. The critical-review stage was especially useful. It pushed the agent to look back at its own results, find weak spots, improve the figures, and strengthen the technical explanations. Without that stage, I probably would have accepted the first report, even though it could have been much better.

Technically, the assignment deepened my understanding of Bayesian classification and uncertainty estimation. By studying the last-layer Laplace approximation, I saw how priors, likelihoods, and Bayes\textquoteright{} theorem connect to posterior inference. I also got a clearer picture of maximum a posteriori estimation, Hessian matrices, Gaussian posterior approximations, and Monte Carlo prediction. One point that I initially misunderstood was that computing an exact conditional Hessian does not make the posterior itself exact. The Laplace approximation is still a local Gaussian approximation, and it does not fully capture uncertainty in the learned feature representation.

Another lesson was that a more sophisticated algorithm does not always produce better results. The Laplace classifier reached about 97.69\% recognition accuracy, but its negative log-likelihood was higher than that of the MAP classifier. That result was humbling. It showed me that accuracy alone is not enough, especially for probabilistic models. Negative log-likelihood, Brier score, and calibration error tell us something important about whether the predicted probabilities can be trusted. I also learned to value controlled comparisons, repeated experiments, validation-based model selection, and careful error analysis. Even negative experimental findings can be useful if they are supported by evidence and explained clearly.

Overall, this assignment changed how I think about AI in programming and research. The quality of AI-generated work depends not only on how strong the model is, but also on how well the user organizes instructions, provides relevant skills, and defines verification standards. An AI agent can automate a lot of implementation, experimentation, and documentation, but I still need to understand the underlying concepts, question the methodology, and interpret the results. I now see AI as a powerful assistant, not a replacement for human judgment. Its usefulness depends on clear direction, reliable verification, and critical thinking.
